% ============================================================
%  Příprava na maturitu – Mechatronika
% ============================================================

\documentclass[a4paper,10pt]{article}

% --- Balíčky ---
\usepackage[T1]{fontenc}
\usepackage[utf8]{inputenc}
\usepackage[czech]{babel}
\usepackage[left=2.4cm, right=1.8cm, top=1.8cm, bottom=1.8cm]{geometry}
\usepackage{lmodern}
\usepackage{microtype}
\usepackage{parskip}
\usepackage{titlesec}
\usepackage{enumitem}
\usepackage{xcolor}
\usepackage{hyperref}
\usepackage{tabularx}
\usepackage{array}
\usepackage{tikz}
\usepackage{eso-pic}
\usepackage{mdframed}
\usepackage{amsmath}
\usepackage{amssymb}

% --- Barvy ---
\definecolor{accent}{HTML}{03254E}
\definecolor{stripeblue}{HTML}{568EA3}
\definecolor{medgray}{HTML}{888888}
\definecolor{lightblue}{HTML}{EAF2F8}
\definecolor{lightyellow}{HTML}{FDFBE4}
\definecolor{lightgreen}{HTML}{EAF7EA}

% --- Levý pruh ---
\AddToShipoutPictureBG{
  \begin{tikzpicture}[remember picture, overlay]
    \fill[stripeblue]
      (current page.north west)
      rectangle
      ([xshift=10mm] current page.south west);
    \fill[accent!60]
      ([xshift=10mm] current page.north west)
      rectangle
      ([xshift=12mm] current page.south west);
  \end{tikzpicture}
}

\titleformat{\section}
  {\large\bfseries\color{accent}}
  {\thesection.}{0.5em}{}[\color{accent}\titlerule]

\titlespacing{\section}{0pt}{16pt}{6pt}

% --- Boxy ---
\newmdenv[backgroundcolor=lightblue,linecolor=stripeblue]{pojmybox}
\newmdenv[backgroundcolor=lightyellow,linecolor=accent!40]{vysvetlenibox}
\newmdenv[backgroundcolor=lightgreen,linecolor=accent!30]{vzorcebox}

% ============================================================
\begin{document}

\section{Metody technické diagnostiky}

\begin{pojmybox}
\textbf{Klíčové pojmy:}
prediktivní údržba, vibrodiagnostika, termodiagnostika,
tribodiagnostika, akustická emise, NDT,
spolehlivost, MTBF, MTTR, Weibullovo rozdělení
\end{pojmybox}

\begin{vysvetlenibox}

% ------------------------------------------------------------
\textbf{Strategie údržby}

\begin{center}
\begin{tabularx}{0.95\linewidth}{l X X}
Strategie & Popis & Náklady \\
\hline
Reaktivní & oprava až po poruše & velmi vysoké \\
Preventivní & pravidelná výměna & střední \\
Prediktivní & dle stavu zařízení & nejnižší \\
\end{tabularx}
\end{center}

\textbf{Shrnutí:}
\begin{itemize}
\item cílem je přejít z reaktivní → prediktivní údržbu
\item využívají se senzory + analýza dat
\end{itemize}

% ------------------------------------------------------------
\textbf{Ukazatele spolehlivosti}

\begin{vzorcebox}
$MTBF = \frac{t_{provoz}}{N_{poruch}}$

$MTTR = \frac{t_{opravy}}{N_{oprav}}$

$A = \frac{MTBF}{MTBF + MTTR}$
\end{vzorcebox}

\textbf{Význam:}
\begin{itemize}
\item MTBF = jak často se to rozbije
\item MTTR = jak dlouho to opravuješ
\item dostupnost = jak moc je stroj „k dispozici“
\end{itemize}

% ------------------------------------------------------------
\textbf{Vibrodiagnostika}

\begin{itemize}
\item měření vibrací (akcelerometr)
\item analýza frekvence (FFT)
\item každá porucha má typický signál
\end{itemize}

\begin{center}
\begin{tabularx}{0.95\linewidth}{l X}
Závada & Projev \\
\hline
Nevyváženost & 1× otáčky \\
Nesouosost & 2× otáčky \\
Ložiska & specifické frekvence (BPFO, BPFI) \\
Ozubená kola & násobky zubů \\
\end{tabularx}
\end{center}

% ------------------------------------------------------------
\textbf{Termodiagnostika}

\begin{itemize}
\item měření teploty (IR kamera)
\item odhaluje:
\begin{itemize}
\item přehřáté ložiska
\item špatné kontakty (elektrika)
\item tření
\end{itemize}
\end{itemize}

% ------------------------------------------------------------
\textbf{Tribodiagnostika}

Analýza oleje:

\begin{itemize}
\item obsah kovových částic → opotřebení
\item viskozita → degradace oleje
\item voda → porucha těsnění
\end{itemize}

% ------------------------------------------------------------
\textbf{NDT metody (nedestruktivní testování)}

\begin{itemize}
\item RT (rentgen) → vnitřní vady
\item UT (ultrazvuk) → trhliny, tloušťka
\item MT → feromagnetické materiály
\item PT → povrchové trhliny
\end{itemize}

% ------------------------------------------------------------
\textbf{Weibullovo rozdělení}

Používá se pro modelování poruch:

\begin{itemize}
\item rané poruchy (výrobní vady)
\item náhodné poruchy
\item opotřebení (stárnutí)
\end{itemize}

\textbf{Tzv. vanová křivka (bathtub curve):}
\begin{itemize}
\item začátek → vysoká poruchovost
\item střed → stabilní provoz
\item konec → opotřebení
\end{itemize}

% ------------------------------------------------------------
\textbf{Shrnutí}

\begin{itemize}
\item cílem diagnostiky je zabránit poruchám
\item využívá se měření (vibrace, teplota, olej)
\item moderní trend = prediktivní údržba + data
\end{itemize}

\end{vysvetlenibox}

\end{document}
